\chapter{Introduzione}

\section{Obiettivo del documento}
Questa documentazione descrive il comportamento della UI avviata da \texttt{ui\_config\_tool.py}.
La struttura segue il flusso reale dell'applicazione: un capitolo per ogni finestra apribile dal menu principale e sottocapitoli per le sottofinestre attivabili nelle varie schermate.

\section{Ambito}
Il documento copre:
\begin{itemize}[leftmargin=1.5em]
  \item finestre principali della UI;
  \item dialog e sottofinestre operative;
  \item relazioni tra configurazione locale, file temporanei e workflow farm.
\end{itemize}

Il documento non copre nel dettaglio l'implementazione interna dei moduli di servizio, se non quando necessario per spiegare input/output utente.

\section{Convenzioni usate}
\begin{itemize}[leftmargin=1.5em]
  \item I nomi dei pulsanti sono riportati come appaiono in interfaccia.
  \item I file JSON sono indicati con percorso relativo alla root del progetto.
  \item Le cartelle operative persistenti e temporanee sono descritte in appendice.
\end{itemize}
